\chapter{What Exactly Is the Self?}
\section{A Student ID Card}
Begin with an object: your student card or identity card.

It carries a photograph, a name, a birth date, perhaps a number. You bring it through the school gates, swipe it at the library, and show it to invigilators comparing the photograph with your face. The whole world agrees: the person this card identifies is you.

Now do something small. At home, find a photograph of yourself at three. Most families have several: a chubby child with gleaming round cheeks and thin hair, stuffing cake into their mouth. Your family says: look, that was you when you were little.

But examine the photograph, then your reflection. Almost every cell in that child's body has since been replaced. The child knew neither multiplication tables nor examinations nor your present best friend. Their likes, fears, and voice scarcely overlap with yours. If that three-year-old suddenly appeared before you, you probably would not recognise them, and they certainly would not recognise you.

What, then, entitles us to call that child ``you''?

The student card answers briskly: a number, a file, successive official records connecting your three-year-old and present selves. But that merely passes the question to the system. On what does the system base its verdict? A shared body? Its materials have changed. Shared memories? You cannot remember what you thought at three. A resemblance? Honestly, there is little.

You use ``I'' every day, perhaps hundreds of times: I am hungry; I think this is unfair; I want to become a doctor. The word comes effortlessly, yet ask what exactly it names and most people suddenly stall. This chapter takes apart a word everyone possesses but almost nobody has inspected inside.

\section{A Birthday in a Hospital Room}
Come into a scene.

An October afternoon, bright with sunshine. A hospital's third floor, a two-bed room. Disinfectant scents the corridor. A nurse pushing a treatment trolley presses a hand-sanitiser dispenser by the door with a soft click.

In the bed beside the window sits Mr Chen, in his seventies. Today is his granddaughter's birthday, and the family has gathered. His son carries a cake box, his daughter-in-law flowers. Sixteen-year-old Xiaoyu still wears her uniform, having come straight from school.

The old man seems cheerful, smiling as they enter. His son sets the cake on the bedside cabinet. ``Dad, Xiaoyu has come to see you.''

Mr Chen turns towards her and smiles politely. ``Thank you, young lady. Whose child are you?''

The room falls silent for a second. Last year Xiaoyu still visited weekly to play chess; Grandpa taught her the Chinese-chess horse's L-shaped move. Now his eyes show no recognition. He is neither pretending nor sulking. Three years ago doctors diagnosed an illness that slowly erases memory: yesterday first, then last year, then farther and farther backwards.

Her mother says softly, ``Dad, this is Xiaoyu, your granddaughter.''

``Oh, oh.'' He nods, still polite. ``So grown up.'' He sounds as though he has met an old colleague's child in the street.

When they blow out the candles, he happily claps and joins the birthday song, though he misses some words. Xiaoyu later tells her mother that what hurt most was not his failure to recognise her. It was that this failure did not seem to make ``Grandpa'' disappear. He still secretly gives the cake's strawberries to children, sings high notes out of tune, and raises his left eyebrow higher than his right when smiling.

Remember this room. It poses an exceptionally difficult philosophical problem: is someone who has lost most memories still the original person? Of course, says the family; that is why they visit weekly. But if memory does not ground ``he is still himself,'' what does? And if one day the raised eyebrow and shared strawberries disappear too, is the person lying there still the same Grandpa Chen? On which day would he cease to be?

\section{Cornering the Question}
Do not demand an answer yet. First make the conflict clear.

Several answers to ``What am I?'' coexist in your mind. Usually they get along; in a hospital room they fight.

\textbf{Answer One: I = this body.} It has existed continuously since birth: born in hospital, attending primary school, now holding this book. Wherever the body is, I am. This answer is reassuringly concrete. Police investigations, emergency treatment, and examination identity checks rely upon it.

\textbf{Answer Two: I = my memories.} I am myself because I remember where I came from: first-day nerves at primary school, last summer's holiday, promises to friends. Forget everything and you lose yourself. This explains pride or shame about childhood: those were my experiences.

\textbf{Answer Three: I = the person in other people's eyes.} A parents' son, a teacher's student, a friend's closest companion. Relationships form a net, and I am its knot. This grounds the family's conviction that the man in bed remains Grandpa Chen.

Usually all three identify one person, so no choice is needed. The hospital tears them apart. The body remains---Answer One says yes. Memory is damaged---Answer Two says not entirely. Relationships continue---Answer Three says yes, always. Three answers produce three verdicts, each reasonable.

This is not invented wordplay. Real questions depend on it. Must someone with amnesia answer for earlier deeds? Has someone whose character changes radically ``broken the former self's promises''? Where should law, morality, and affection draw the boundary of the same person?

There is a deeper layer. All three answers ask how to judge whether present and past selves are the same---a question of \textbf{continuity}. Solve it, and something more basic remains: is any ``thing'' travelling along that continuous line? We call a train from Beijing to Guangzhou the same train because a metal vehicle travels the route. In the train called ``I,'' what is that enduring metal vehicle? A soul? A self? Or is there no vehicle, only successive scenery, with ``passenger'' merely a story we invent?

Some of humanity's earliest thinkers disputed this long ago. Their answers clash, and no winner has emerged. Our next stop is India over two millennia ago.

\section[Two Opposing Answers in India]{Beside the Ganges and beneath the Bodhi Tree: Two Opposing Answers}
Turn back twenty-five to thirty centuries, to India. Two of history's most opposed answers to ``What am I?'' arose here, on almost the same land and only centuries apart.

\textbf{First answer: A true self lies within you, unlike what you imagine and vastly larger than you think.}

In ancient texts later called the \textit{Upanishads}, teachers repeatedly discuss two terms. \textit{Atman} roughly means a person's deepest self: not name or appearance but what sees, hears, and experiences beneath all thoughts, memories, and moods. \textit{Brahman} names the ultimate reality underlying the universe. Their startling claim is that these are the same. You think ``I'' is a small tenant locked inside your body; dig to the deepest level and you reach the foundation of the universe itself.

The \textit{Chandogya Upanishad} offers a famous lesson. A father asks his son to dissolve salt in water. Next day the salt is invisible, yet every part tastes salty. The unseen self is like salt: invisible but everywhere. Throughout the teaching, the father repeats ``Tat tvam asi,'' broadly ``You are that.'' You, the questioning child before him, are in your deepest being that ultimate reality. On this tradition's account, suffering and confusion arise from mistaken identity: taking body, reputation, and emotions for the self while forgetting the true self beneath. Practice means finding it again.

\textbf{Second answer: Stop searching. Dig to the bottom and find nothing, because no fixed self was ever there.}

This is a central Buddhist claim: non-self, Sanskrit \textit{anatta}, Pali \textit{anatta}. It does not mean you do not exist. You are plainly reading, becoming hungry and sleepy. It means that no part of you is \textbf{unchanging and independently existent} enough to deserve the title ``true self.''

Take it apart. Is the body self? It grows, falls ill, ages, and constantly replaces material; if it were self, would I not become someone else every minute? Feelings? Pleasure and boredom take turns. Memory? Grandpa Chen shows its unreliability. Thoughts? Five years ago you may have fiercely rejected what you now firmly believe. Buddhism groups a person's components into five heaps---body, feeling, perception, habits, consciousness---then observes that each flows and none commands everything. Adding them cannot produce an eternal, unchanging self. ``I'' is a useful label attached to this flow, but no iron core lies beneath it.

How far apart are these answers? One says: search deeply and the true self exceeds imagination. The other: search deeply and discover that no such true self exists. They grew from the same soil, read each other's writings, and argued for over a millennium. This is no uniformly voiced ``Eastern wisdom.'' Its first great internal debate concerned the self.

Now perform an important exercise: state the full case for whichever side you trust less.

Suppose non-self appeals more, while soul and true self seem antiquated. Stand with the Upanishads and strengthen their case. First, they address an undeniable experience: however bodies and thoughts change, something seems continuously present, experiencing it all. The perspective looking at cake at three feels like the perspective looking at the photograph now. That experience needs explanation; atman supplies one. Second, it deeply grounds morality. If everything shares one ultimate reality, harming another literally harms yourself. Compassion needs no justification; it is a fact. Third, it accounts for death. The body disperses, but the deepest self is unborn and undying. For someone afraid of disappearing, this is support, not superstition.

Conversely, if you favour a true self, give Buddhism its full argument. First, it is honest: no unchanging thing has ever been found in body or brain. Insisting on a core may express a wish, not a discovery. Second, it identifies the practical damage of attachment to ``I.'' The more tightly people grasp face, possessions, and inviolable opinions, the more they suffer. Much conflict and anxiety arises from this grasping. Third, it does not declare everything unimportant. Precisely because no fixed self exists, change is possible; your present self is not pinned down by the past.

Correct a common misreading as we pass. People hear non-self as ``Humans are illusory, so actions do not matter and responsibility disappears.'' Buddhism itself fiercely rejects that conclusion. It strongly emphasises responsibility: without a fixed self, each present choice \textbf{makes} the next moment's you, and nobody can cover for you. Hearing an anti-fatalist doctrine as permission to give up is this passage's easiest mistake, ancient or modern.

\section[Thinking Bridge: Three Identity Ledgers]{Thinking Bridge: Auditing the Three Ledgers of the Same Self}
We now need a portable tool, or ``What am I?'' remains a mystical tangle.

The tool is simple. When asked whether someone remains their former self, or whether you will still be you in ten years, stop arguing vaguely. Divide ``the same'' into three independent ledgers and inspect each.

\textbf{First: the bodily ledger.} Has this body existed continuously, from yesterday to today, from age three to now, as an unbroken physical process? Fingerprints, records, scars, and relatives' identification make this comparatively easy to audit. Legal systems mainly operate through it.

\textbf{Second: the memory ledger.} Can the person connect inwardly with the past? Do they remember an event as something they did, rather than something reported to them? Grandpa Chen's bodily ledger is intact, but his memory ledger has large deficits.

\textbf{Third: the character-and-relationships ledger.} Do temperament, habits, and values continue? Does the person still occupy the same place in others' lives: the grandfather who shares strawberries, Xiaoyu's grandfather, his son's father?

The tool helps because disputants often use different ledgers while imagining they calculate the same thing. Family saying ``Grandpa is still Grandpa'' use the first and third; observers saying ``He is no longer the same person'' use the second. Both ledgers are real, but their conclusions conflict. Distinguish them and at least the argument reaches its actual point.

More importantly, the tool brings an unsettling discovery: \textbf{none of the three ledgers is indestructible}. The bodily ledger is steadiest, yet cells and appearances change, and severe injury, cosmetic surgery, or transplantation can blur it. Memory seems most inward but is least reliable. Psychologists have long shown that remembering reconstructs rather than replays. People sincerely remember events that never happened. In experiments, adults shown a doctored childhood hot-air-balloon photograph, with suggestive prompting, often recalled details after several interviews. Your trusted ``I remember'' may be a construction crew, not an archive. Relationships live in other minds; misunderstanding, rupture, or forgetting changes their ledger too.

If the self is property, all three account books are handwritten, subject to alteration and loss. Keep this unsettling fact in mind; every later section needs it.

\section{Reading a Two-Thousand-Year-Old Dialogue}
Now read a short primary passage. Non-self is not merely a label later readers imposed on ancient people. They explained it themselves, more elegantly than textbooks.

The \textit{Discourse of the Monk Nagasena} records a vividly plausible exchange. Milinda, a king deeply influenced by Greek culture, questions the monk Nagasena. What are you called? People call me Nagasena, he replies, but that is only a name; no fixed, unchanging person is here.

The king objects: without a person, who eats, observes discipline, and experiences the consequences of actions?

Rather than argue directly, Nagasena asks whether the king came by carriage or on foot. By carriage. Nagasena then proceeds, broadly:

\begin{quote}
Is the axle the carriage? No. Are the wheels the carriage? No. The body? The shafts? The reins? None. Is the sum of all these parts the carriage? No: piled parts remain a pile. Then where is the carriage existing separately from its parts? Nowhere to be found.
\end{quote}
The king concedes that no part is the carriage, nor is there a carriage itself hidden behind the parts. ``Carriage'' is the name given to their assembled arrangement.

People are alike, says Nagasena. Beneath his name are hair, skin, bones, feelings, thoughts, and consciousness. Examine each: none is Nagasena. Add them: no Nagasena appears. Apart from them, none exists either. The name is useful, as ``carriage'' is useful, but names a functioning combination, not an entity.

The dialogue is classic because it uses no appeal to authority---``Scripture says so''---but an inspection anyone can follow. You claim a self exists? Point it out. Every candidate exposes a problem. Non-self becomes a hands-on task rather than a declaration of faith.

Notice its boundary too. It dismantles a ``carriage spirit behind the parts,'' but leaves another question: if no carriage itself exists, why call these parts by one name for five or ten years? It removes substance, not continuity, precisely what the hospital room needs most. Carry that question onward.

\section[Soul, Relationships, Individual]{Soul, Relationships, Individual: Three Traditions and Three Selves}
The same facts---bodies die, memories fail, relationships change---produce different civilisational blueprints of self. Compare three influential ones. They are not three expressions of one idea but different verdicts on what the self is, each with reasons and costs.

\textbf{Blueprint One: The Greek soul---an immortal inner component.}

In Plato's \textit{Phaedo}, Socrates calmly discusses death before execution. Broadly, his central claim is that the soul is the true person, the body merely temporary clothing. Because the soul thinks and knows eternal truths, it too is immortal. Death ends not you but the soul's wearing of its clothes. In the \textit{Phaedrus}, Plato offers another famous image: the soul is a chariot, reason its driver, controlling two horses, one noble (spirit), one unruly (desire). Being yourself means learning to drive.

Its strengths: a stable core explains the feeling of a master within the body; morality gains a support, since controlling desire is the driver's duty; death receives an explanation. Its costs: separating wearer from clothing can devalue body, desire, and earthly life. Moreover, nobody in over two millennia has pointed out what or where this core is. It requires belief, not inspection. Nagasena's procedure directly challenges such blueprints.

\textbf{Blueprint Two: The Confucian relational self---a knot formed in a network.}

Pre-Qin Confucianism follows a different route. It scarcely asks what an isolated self is, because such a thing hardly exists in its view. In ``Yan Yuan'' in the \textit{Analects}, Confucius answers a question about government in eight Chinese characters:

\begin{quote}
Let rulers be rulers, ministers ministers, fathers fathers, and sons sons.
\end{quote}
In essence: each should act as their role requires. Often treated as a hierarchical slogan, this conceals a claim about self: relationships and actions define a person. You are a father not because you contain a father-core, but because you nurture and take responsibility. Fully inhabiting roles establishes the person. The four-character phrase \textit{ren zhe, ren ye}, ``benevolence is humanity,'' in the \textit{Book of Rites}'s ``Doctrine of the Mean'' goes further. Humanity's fundamental quality, benevolence, is becoming human itself: responsive connection with others, the tightening in your heart at their pain. On this blueprint, an isolated self does not first exist and then enter relationships; relationships gradually weave you. A parents' son, a friend's friend, a teacher's student: the sum of these identities is you.

Its strengths: it explains what many Western blueprints cannot, including why the family insists Grandpa Chen remains himself. In the relational ledger he has never left. It also explains why Chinese speakers discussing ``I'' often instinctively include ``our family'' or ``our workplace.'' Its costs: the network can swallow the individual. Define identity wholly by roles and someone who rejects being the obedient son to become themselves struggles to find footing. Networks can oppress; ``Behave as you should'' may educate or imprison.

\textbf{Blueprint Three: The modern individual---an atom bearing its own rights.}

This younger blueprint developed in Europe over roughly the past three or four centuries and now spreads through globalised schools, laws, and phone apps. Each person is first an independent \textbf{individual}: self-owning, choosing by personal will, carrying rights and dignity. Relationships are \textbf{chosen}, not what generates you. You may leave your birthplace, change faith, redefine yourself. ``Be your authentic self'' becomes a supreme instruction. Every app registration assumes it: one account, one person, one ``I agree.''

Its strengths: it shields ordinary people, including the unwilling obedient son. Before any role, you are yourself; nobody may use you merely as a tool. Opposition to arranged marriage, protection of privacy, and recognition of minors' independent personhood rely upon it. Its costs: atoms are lonely. Treat the self as an independent unit with its own manual and human ties can look like burdens. When someone ages, loses memory, or loses autonomy, their value suspiciously shrinks within this blueprint. That is our hospital problem: if autonomous choice wholly supports the self, what remains of Grandpa Chen without that capacity? The family's ``He is still here'' plainly uses another blueprint.

Together they reveal something subtle: each solves real problems where it works best and creates others where it works least well. It is also wrong to assign one blueprint to an entire civilisation throughout history. Every tradition argues internally. Confucianism contains voices emphasising personal attainment; Western thought has long criticised atomistic individuality. Blueprints are toolboxes, not household registration documents.

\section[Further Challenge: Thought Experiments]{Further Challenge: Dismantling the Self in a Thought Laboratory}
When theories deadlock, philosophers have a cheap method: \textbf{thought experiments}. Construct an extreme situation mentally and see whether intuition overturns. Try operating three of the discipline's best-known machines.

\textbf{Experiment One: Change bodies.} In \textit{An Essay Concerning Human Understanding}, seventeenth-century English philosopher Locke imagines, broadly, a prince's complete memories entering a cobbler's body overnight. The cobbler wakes remembering palace life, commands issued, people feared, but nothing about the calluses on his hands. Who wakes: prince or cobbler?

Locke answers clearly: the prince. Personal identity rests on continuity of consciousness and memory, not body. Notice the explosive implication: memory, not body, becomes the main ledger. If Locke is right, the hospital verdict changes. With large memory deficits, Grandpa Chen's person really is departing. Do you accept that? Before answering, push further. If memory can relocate completely, the self resembles a copyable file. But can a file be copied twice? That leads to the next machine.

\textbf{Experiment Two: Lose memory.} No fantasy is needed; this machine operates daily in hospital rooms we have already visited. Sharpen it with a real medical case type. Some severely brain-injured patients cannot form new memories. Memory stops in a particular year, and every subsequent day is new. Nurses reintroduce themselves daily; patients sincerely say ``Nice to meet you'' each time. Over twenty years, has the person lived one life or one day thousands of times? Externally the bodily ledger continues; internally memory breaks into countless pieces. If self lives in memory, who has lived those twenty years? If you cannot bear to say the person is gone, you are quietly using another ledger. Identify which.

\textbf{Experiment Three: Make a copy.} Twentieth-century philosopher Derek Parfit made this machine famous. Imagine a future transporter: enter a departure chamber, have your entire body scanned, send the data to Mars, and reconstruct an atom-for-atom counterpart from new material in a receiving chamber, complete with memories and character. Meanwhile the Earth original is dismantled. Cheap and fast; everyone travels for business this way. Is the person leaving the Mars chamber you?

Intuition divides. Yes: the arrival has all your memories, first thinks ``I have reached Mars,'' loves those you love, and remembers your passwords. Under Locke's memory ledger, this is you. Why object to new bodily material when cells already change? No is equally forceful: dismantling on Earth is destruction, scanning plus killing, not travel. You were merely read, then disposed of.

Turn the machine to its harshest setting. An upgrade means you are \textbf{not} dismantled on Earth. You dust off your clothes and leave the chamber while a counterpart still appears on Mars. Two of you now share identical memories, diverging from this second. Which is you? ``Both'' removes uniqueness: next month they form separate romances and debts; how do we account for those? ``The Mars one is only a copy''---why? He too wakes into complete memories and feels bitterly wronged by being treated as counterfeit. Now no ledger settles the case: both bodily ledgers and both memory ledgers qualify.

These machines share an uncomfortable but important result: \textbf{our intuition of self has been pampered by an everyday world where there is only one and no branching}. Permit splitting, copying, or interruptions, and intuitions crash. Parfit drew a radical conclusion, broadly: if thought experiments can leave ``Is that me?'' unanswered, perhaps identity matters less than we assume. What matters is not whether the future person is me, but how many psychological connections link their experiences with my present. The self resembles a wave: not a special droplet named ``I,'' but the movement connecting one wave with another. You may disagree, as many philosophers do. Still, he transforms a yes-or-no question into one of degree, a move rehearsed two millennia earlier in Nagasena's dialogue.

\section{Today's Self Lives in a Phone}
The ancient question now reappears daily in your pocket.

First, doubles. At least two versions of you run simultaneously: one in class, one in an account. The account has an avatar, biography, and posting history through which strangers describe you. Does it count as you? Memories and choices generate it, partly satisfying the memory and character ledgers, but it lives on servers, open to screenshots, reposting, and quotation out of context, with no bodily ledger. Many know the strangeness of someone defining them by a sentence posted three years ago: ``This is who you are.'' Your present inward self protests, but the textual version never ages or retracts. Online it outlasts your body and sometimes speaks louder. A three-year-old photograph feels alien; so does a three-year-old post. This time the whole world can assign that stranger back to you.

Second, copies. Programs trained on someone's posts and voice samples can already imitate their conversation convincingly. If one becomes close enough---remembering your jokes, using your catchphrases, comforting as you do---what is its relationship to you? Most instinctively answer that it is only a program. Audit that refusal. Body: it lacks mine? Locke asks when body became the main ledger. Memory: it never actually experienced those events? Your childhood recollections may not all be actual experiences either. Defending ``That is not me'' proves harder than expected, not because it really is you, but because these ledgers were never designed for such objects.

Third, those who have left. Accounts, chats, and voice messages give a deceased person's presence a new intermediate state. The relational ledger keeps operating: friends leave birthday messages, family view photographs. Body and memory cease updating. The Confucian blueprint can name such mourning: while relationships and rites continue, the knot remains in the net. The modern-individual blueprint urges letting go because the person has gone. These blueprints argue in countless living rooms without anyone naming them.

Finally, responsibility. Someone posts hurtful words at fifteen; they resurface at twenty-five. Is this still the same person? Forgive or hold accountable? Body says the same, character may say radically changed, memory depends on whether the person still acknowledges it. Society disputes cancellation and forgiveness, growth and reckoning. Beneath it runs this old machine: which ledger settles personal identity?

\section[Your Turn: Leaving the Transporter]{Your Turn: After Leaving the Departure Chamber}
Return to the transporter and add one final setting before answering.

Technology advances again. You may now travel to Mars or send a duplicate while remaining on Earth. An excellent opportunity awaits on Mars, a project you have long dreamed of, but an ill relative needs you at home. You decide: send a counterpart to pursue the dream and stay yourself. Ten years pass. The Mars person becomes who you once hoped to be, writing home monthly in your exact tone. On Earth you live another life.

Now answer, necessarily with reasons:

\textbf{Does the person on Mars owe this family anything?} On departing, if that counts as his departure, he promised to care for your dreams; you promised to care for relatives on behalf of both. Do those promises still bind each respective person? Suppose he says, ``The shared pre-departure self made the promise. That self no longer exists. We are now two people owing each other nothing.'' Can you refute him? With which ledger?

If that sounds like evading a debt, notice that everyone on Earth uses the logic daily: ``That was the old me,'' ``People change.'' Only degree differs. How much responsibility does a person bear for a former self? When is ``I have changed'' growth, and when is it debt avoidance?

Conversely, if you think he cannot escape, be equally careful. That means your self ten years hence must answer fully for every word and wish you write now, including an impulsive late-night message.

There is no standard answer. One thing is certain: since birth you have accepted bills left by a former self you never meet, while issuing new ones to a future self you never will. Every map in this chapter---Ganges and Bodhi-tree disputes, Nagasena's carriage, three civilisational blueprints, three thought-experiment machines---aims to make you inspect this account you never signed yet continually honour, then ask:

Knowing all this, how will you sign?
